\section{Numerical contract and implementation coverage}
\label{sec:supp-contract}

\subsection{What is bound and what is checked}
The reference binds the producer as well as the consumer. In the CUDA deployment, the original prompt producer uses physical microbatch one and requests its original final output. Candidate prefill can regroup known rows when the resulting requested vector and active cache preserve that producer. The controlled SmolLM2 Vulkan append study retains its original prompt calls of up to 512 tokens and physical microbatch 128. Equal emitted-token endpoints alone do not make states produced by different schedules interchangeable.

The dense append check in \texttt{phase\_append\_requirements} requires one stream, no sliding window or shifted cells, an occupied prefix whose end and physical head equal the append position, and sufficient capacity. Existing cells must have their own logical positions and include sequence zero. The native padding is at least 256 entries. The executor additionally binds F16 K/V and a context owned by its single-sequence frontend. The sequence-membership check is not an exclusivity test for externally mutated, shared caches.

The phase executor accepts all-row or last-row output requests. With zero-based row indices these modes instantiate the main paper's set $O$ as $\{0,\ldots,q-1\}$ or $\{q-1\}$, respectively. The original released Q8 frontend reserves four output rows; its all-output gate rejects wider requests before a native call, although last-only prompt calls can be wider. S3 uses a separate 32-output reserve. An eight-output diagnostic profile likewise has its own consumer and qualification. The planner's more general requested-row flags do not expose an arbitrary sparse-output API through the evaluated frontend.

Local attention requires causal, non-Flash execution, consecutive positions, one sequence, a complete prefix-extent vector, singleton K/V streams, and no attention bias, sinks or MLA path. Its descriptor includes Q/K/V types, head dimensions and counts, and V transposition. Equal-descriptor runs are tiled to at most four query columns in the admitted layout, with K, V and mask views restricted to the required extent. A noncontiguous mask is materialized before use. The planner can check the row partition; equality of the selected kernels' arithmetic remains a backend condition. Several architecture and layout requirements are assertions rather than recoverable admission statuses.

\subsection{Projection checks and reuse}
The Q8 packed branch accepts F32 activations with Q8\_0 weights and query width 5--32. Actual multi-group packing additionally requires contiguous activations, singleton higher weight dimensions and at least eight queries. It maps complete groups to $[K,4,\lfloor q/4\rfloor]$, reshapes the outputs, and handles a remaining tail with an ordinary operation of at most four columns. Widths five through seven therefore enter this branch without constructing multiple packed groups.

The same reference binding must cover the smaller widths created by a physical cut. In the four-query example, \Simple{} sends three and one rows through the graph, while \Local{} sends four. In the eight-row prefill illustration, the physical boundary is after the seventh row; local attention uses $4+3+1$ tiles and projections use $4+4$ groups. The figure shows the grouping construction, not a newly measured workload. Preserving the full-group shape does not prove equivalence for these tails, their quantizers or the surrounding decoder operations.

The marked backend-loop4 control checks contiguous Q8\_0/F32/F32 tensors, dimensions, offset bounds and absence of fused arithmetic or indexed matrix multiplication. Its CUDA dispatcher requires actual buffers on the selected device, NVIDIA warp size 32, safe padding, and MMVQ selection for every original one-to-four-column tile. These checks precede that operator's quantizer and kernel, but occur inside a decode call; earlier graph operations can already have executed.

Packed projections have no loop4 marker. Their ordinary dispatcher can select other kernels, including a cuBLAS branch for unsuitable padded views. Falling back to an ordinary graph path is not, by itself, a reference-equivalence guarantee. The mixed Q4\_K/Q6\_K experiment has its own one-column control and is outside these Q8 marker checks.

The active graph-reuse checks compare the policy, full prefix-extent vector, head-placement key, and normal graph topology and output count. Requested row indices are refreshed as graph inputs. A structural identity helper exists in the planner header but is not the live reuse key. Accordingly, neither a helper's presence nor a cache-type tag supplies a general run-time certificate for every numerical choice.

\subsection{Output lifetime, observation and failures}
After synchronization, each requested vector is copied to storage indexed by its logical row. A later group can then reuse native scratch buffers without destroying an earlier accepted frontier. The implementation checks that a requested output exists and that the physical append endpoint is correct; it does not compare reference KV bytes online. Greedy sampling checks all vector entries for finiteness and retains the first vocabulary index on ties.

The speculative loop materializes an already emitted seed and the draft proposals, samples every requested target row, and then processes acceptance in order. With $b$ initially materialized rows and $r$ accepted proposals, it trims target and draft to length $b+1+r$. A correction or bonus token is emitted after that accepted prefix and remains unmaterialized until the next decode. On a cap or end-of-generation stop, the reported stream can end in this pending state; no extra decode is introduced solely to make KV length equal emitted-stream length.

The original capture observer selects attention and specially named nodes; its filter does not include the packed projection names. Capture is disabled during clean timing. The original state-capture call records initialization and the first commit. Each other profile must be assessed from its own capture implementation and qualification records. A native sequence export also requires extraction of active K/V before it represents the reference relation in the main paper.

The target decode loop poisons its owner on exceptions from native groups and postchecks. Trim, reset, draft, sample and diagnostic failures lie outside that catch; the command-line entry exits on failure. Assertions terminate directly. The resulting process boundary prevents another request in that worker, but does not establish recovery for a library caller that catches an exception and continues.

\subsection{Diagnostics that distinguish consumer contracts}
Two CPU diagnostics demonstrate why output and state requirements must be declared together. They use a synthetic two-layer F32 model, F16 KV, a vocabulary of 128, one CPU thread and a context capacity of 1024. They are mechanism diagnostics, not GPU qualification or performance measurements~\cite{records}.

First, a two-row call and two successive one-row calls start from identical state at history length 248. Although both operation argmax values and the next forced-continuation argmax agree, 242 of 256 output float bit patterns differ. Their final serialized active-sequence states differ in 43 bytes, and the continuation differs in 127 of 128 float bit patterns. The diagnostic does not localize the cause. It demonstrates that token agreement can coexist with different full vectors, state and subsequent computation.

Second, an eight-row operation is compared with a predefined four-plus-four reference at history lengths 248 and 252. All eight shared outputs, active state, geometry and one forced continuation agree. A restricted downstream consumer instead requires the first four vectors, retains five materialized rows, and then consumes fixed token 37. This comparison uses a four-plus-one reference and trims the wider execution to the same five-row frontier. Extra vectors are captured for diagnosis, but the fifth vector is excluded from the downstream contract before execution. It differs in 117 and 128 entries at the two histories, respectively, despite unchanged argmax; the required four outputs, retained state and forced consumer agree. Thus this restricted consumer passes while one requiring the fifth vector does not. These are explicit diagnostic schedules, not an autonomous scalar deployment reference. They establish neither natural speculative acceptance behavior nor a general qualification of an eight-output production verifier.

\subsection{Captured arithmetic paths and attention repairs}
\label{sec:supp-historical-mechanisms}
The historical RTX width study captures Qwen3-0.6B-Q8\_0 and Qwen3-8B-Q8\_0 on an RTX 5070 Ti Laptop GPU. A deterministic 32-token fixture supplies real model activations. Each model contributes layer-0 Q/K/V, attention-output, FFN gate/up/down and vocabulary projections: 16 descriptors, rather than 16 independent workloads. Replays hold each captured weight and F32 activation prefix fixed. The scalar route evaluates and concatenates single columns; the native route uses a wide two-dimensional matrix; the grouped route uses $[K,4,q/4]$ and restores column order. Recorded complete-output digests give Table~\ref{tab:history-width}~\cite{records}.

\begin{table}[htbp]
\centering\small
\caption{Native width can change arithmetic; outer grouping preserves the sampled outputs. Equality counts compare the complete operator output with same-operand scalar replay across both models. CUPTI dispatch was profiled for the eight 8B descriptors. A dash means no multi-group packed operation was constructed at that width.}
\label{tab:history-width}
\begin{tabularx}{\textwidth}{@{}rYrr@{}}\toprule
Width & Native dispatch observed on 8B & Native = scalar & Grouped four = scalar\\\midrule
\HistoryWidthRows
\bottomrule\end{tabularx}
\par\smallskip\raggedright\footnotesize
\textsuperscript{a}Stream-K fixup occurs for seven of eight traced classes; the vocabulary projection omits it. Grouped widths 8/16/32 retain row-Q8\_1 and four-warp MMVQ. Internal temporary quantized bytes were not separately captured.
\end{table}

All 16 native width-32 replays reproduce their own live captures. At the first Q projection, both models additionally have matching live scalar/native input bytes and scalar-output replay. Later operator comparisons hold their captured operands fixed; they do not reconstruct a whole scalar decoder. The historical source uses four reduction warps at native widths 1--4 and two at widths 5--8. Thus width eight can alter reduction organization even with the same activation-quantizer entry point.

The Vulkan Llama companion starts from its own stock scalar reference: Meta-Llama-3.1-8B-Instruct-Q8\_0 on B570, a CPU Llama-3.2-1B-Instruct-Q8\_0 draft, capacity 1024, F16 KV and $k=3$. A scoped F32 projection policy repairs the scalar-versus-wide activation path. Endpoints then agree, but the long development input still exposes 99 different committed vectors in each speculative layout. Combining that policy with reference attention groups removes the differences (Table~\ref{tab:history-repair}). At the diagnosed boundary, changing extent also changes V contiguity and the selected matvec kernel. This is a joint attention-path repair.

\begin{table}[htbp]
\centering\small
\caption{Projection and attention repairs are complementary in the historical Vulkan case. ``Match'' refers to emitted-token and stopping endpoints. Vector columns count differing complete vocabulary vectors. The long input is CodeContests/valid/108; each of its two speculative layouts contributes 128 committed decisions. The full audit covers three inputs and five audited arms (R, M\_PQ, M\_QB, C\_PQ, C\_QB), including the R control.}
\label{tab:history-repair}
\begin{tabularx}{\textwidth}{@{}Yp{1.7cm}rr@{}}\toprule
Configuration & Endpoints & Long input, per layout & All committed rows\\\midrule
\HistoryRepairRows
\bottomrule\end{tabularx}
\par\smallskip\raggedright\footnotesize
\textsuperscript{b}The same-binary attention-off run reproduces the projection-only trajectory across all 36 requests and 3,152 saved target/draft vectors. Its committed counts follow that recorded trajectory equality. The dash in the projection-off row means no committed-vector count is used here; its endpoint failures already distinguish it. The two layouts' 99 affected positions account for 198 differing records, not 99 across the entire audit.
\end{table}

\subsection{Producer state changes consumer compatibility}
A historical B570 Qwen3-8B/0.6B Q8 test exports the initial target state into a clean receiving context. P1 produces the prompt one token per call; P4T uses four-token calls with separately defined tails. C1 is ordinary target-only continuation; PQ and QB are the two speculative layouts with $k=1$. Table~\ref{tab:history-producer} groups all eight inputs by endpoint outcome against the same independently executed reference R. PQ and QB agree with one another on endpoints, cycles and complete vectors for all 16 input--producer pairs, yet this agreement does not imply equality with R.

\begin{table}[htbp]
\centering\small
\caption{Producer $\times$ consumer endpoint compatibility at the historical boundary-zero clean-receiver interface. Each column compares with the fixed R endpoint. Grouped rows cover all eight inputs and all 48 producer/consumer cells. ``PQ/QB'' means both layouts have the displayed outcome.}
\label{tab:history-producer}
\begin{tabularx}{\textwidth}{@{}Yrcccc@{}}\toprule
Input class & $n$ & P1 C1 & P4T C1 & P1 PQ/QB & P4T PQ/QB\\\midrule
\HistoryProducerRows
\bottomrule\end{tabularx}
\end{table}

On the two prealgebra inputs, changing the producer changes whether the speculative consumers preserve R while C1 retains the same endpoint. KV serialization and the initial full logits both change on every input; this experiment does not isolate KV alone. These boundary-zero, $k=1$ observations precede the later repairs and belong to a different contract from the capacity-4096 multi-turn implementation. The included matrix retains all individual input identities and producer-conditioned comparisons.
\FloatBarrier
